\begin{question}
\noindent A satellite dish tracking station is represented by the external point $P$, from which two signal paths $PA$ and $PB$ are drawn tangent to a circular exclusion zone with centre $O$, touching the circle at points $A$ and $B$.

The angle $AOB = 148^\circ$, as shown in the diagram.

\begin{center}
\begin{tikzpicture}[scale=1.6]
    % Define the center of the circle
    \coordinate (O) at (0,0);

    % Draw the circle
    \draw (O) circle (2cm);
    \fill[black] (O) circle (2pt) node[below right]{$O$};

    \coordinate (A) at (240:2);
    % Angle AOB = 148 degrees, so B would make 240-148 = 92 with the horizontal
    \coordinate (B) at (92:2);

    % Tangent directions at A and B (perpendicular to OA and OB)
    \coordinate (T1) at ($(A)!1!240:(O)$);
    \coordinate (T2) at ($(B)!1!92:(O)$);

    %Drawing tanents from A and B, slightly longer that the actual tangent length, so that we can ensure an intersection, i.e obtain P. The length of tangent will be tan74*2(radius) = 6.97.. about 7cm
    \coordinate (E1) at ($(A)!7cm!90:(O)$);
    \coordinate (E2) at ($(B)!7cm!-90:(O)$);

    % Find intersection point P
    \path [name path=line1] (A) -- (E1);
    \path [name path=line2] (B) -- (E2);
    \path [name intersections={of=line1 and line2,by={P}}];

    % Draw tangent lines from P
    \draw (A) -- (P) -- (B);

    % Draw radii
    \draw (O) -- (A);
    \draw (O) -- (B);
    \draw (O) -- (P);

    % Labels
    \node[below left] at (A) {$A$};
    \node[above left] at (B) {$B$};
    \node[above] at (P) {$P$};

    % Mark angle at O
    \pic [draw, angle eccentricity=1.4, angle radius = 0.6cm, "$148^\circ$"] {angle = B--O--A};

\end{tikzpicture}
\end{center}

\noindent Calculate the size of angle $APB$.

\vspace{4cm}

\begin{flushright}
    \answerequnit{\angle APB}{$^\circ$}{3}
\end{flushright}
\end{question}